\section{Research Question}
\label{sec:receipts-question}

\subsection{Problem Statement}
\label{sec:receipts-question-problem}

The central problem addressed by the receipts corpus is the absence of a durable, tamper-evident
chain of custody for process-mining claims entering a board-level M\&A transaction. In current
enterprise practice, process-mining outputs are reported as dashboard metrics or spreadsheet
summaries with no cryptographic binding to the event logs from which they were derived, no
binding to the specific discovery or conformance algorithm that produced them, and no
independent replication proof. A financial buyer relying on such outputs cannot verify that
the reported EBITDA impact reflects an actual process-conformance measurement rather than a
manually entered or post-hoc adjusted figure. This problem is not merely an audit inconvenience;
it is a structural defect in the M\&A due diligence chain. The receipts corpus addresses it
directly by requiring every M\&A claim to carry a BLAKE3-hashed event log reference, a
BLAKE3-hashed process model reference, a \texttt{wasm4pm\_query\_result} payload, and an
Ed25519 \texttt{validator\_signature} before the claim may be admitted to the board record.

\subsection{Old-Regime Assumption Challenged}
\label{sec:receipts-question-old-regime}

The old-regime assumption challenged by this research is that \emph{reported results are
trustworthy by virtue of their source system}. Under the old regime, a process-mining vendor
asserts fitness of $0.98$ in a slide deck, and the buyer accepts this assertion because the
vendor is reputable and the tool is commercially licensed. No event log is provided. No process
model hash is supplied. No independent replication is possible. The assertion is taken on faith
because the alternative---manual reconstruction of the entire conformance checking pipeline from
raw transaction logs---is prohibitively expensive.

This assumption is challenged by the thesis that trustworthiness must be \emph{manufactured
into} the result at the point of computation, not asserted after the fact. The manufacturing
approach, implemented through the \texttt{ConformanceAdmissionGate} and the receipt serialization
pipeline using JCS canonical form (RFC 8785), produces results that can be independently
verified by any party with access to the original event log and the published process model hash.
The Ed25519 signature scheme ensures that a post-hoc alteration of any field in the JSON receipt
invalidates the signature, making tampering detectable without access to the signing authority.

\subsection{Hypothesis}
\label{sec:receipts-question-hypothesis}

The primary hypothesis of the receipts corpus is: \emph{A cryptographic receipt chain
employing Ed25519 validator signatures, BLAKE3 content-addressed event log references,
and JCS-canonical JSON serialization is sufficient to make process-intelligence claims
board-admissible in M\&A due diligence contexts, provided that the underlying conformance
authority passes fitness $\geq 0.95$ and precision $\geq 0.90$ thresholds on a replicable
event log.}

A secondary hypothesis is that the manufacturing of wasm4pm modules can itself be receipted
through generation receipts that bind module version, Blake3 hash, test pass count, and
cargo compilation status, such that the module's manufacturing pipeline is as auditable as
the claims it produces.

Both hypotheses are evaluated against concrete evidence: five signed JSON receipts with
fitness scores ranging from $0.975$ to $1.0$ across five M\&A claim domains, and four
wasm4pm generation receipts documenting successful compilation and test passage rates of
$8/8$ for the conformance authority and $5/5$ for the Blue River Dam orchestrator.

\subsection{Success Criteria}
\label{sec:receipts-question-success}

Success for the receipts corpus is defined by four criteria that must all be met simultaneously.
First, all five M\&A claim receipts must carry Ed25519 \texttt{validator\_signature} fields and
BLAKE3 hash fields, with fitness $\geq 0.95$ and precision $\geq 0.90$ on each. Second, the
conformance authority module that underwrites those receipts must pass all unit tests and compile
without errors under \texttt{cargo check}. Third, the \texttt{RECEIPT\_REGISTRY.md} must
enumerate all seven canonical receipt types with named witnesses. Fourth, the
\texttt{ma\_deck\_rendering\_authority\_assessment.md} must confirm the board presentation
pipeline is ready for execution. All four criteria are met as documented in
\texttt{checkpoints/ALIVE\_GATE\_ASSESSMENT.md} with 12/12 gate criteria satisfied and the
\texttt{PROCESS\_INTELLIGENCE\_ALIVE\_001} checkpoint sealed.
